% The document class supplies options to control rendering of some standard
% features in the result.  The goal is for uniform style, so some attention 
% to detail is *vital* with all fields.  Each field (i.e., text inside the
% curly braces below, so the MEng text inside {MEng} for instance) should 
% take into account the following:
%
% - author name       should be formatted as "FirstName LastName"
%   (not "Initial LastName" for example),
% - supervisor name   should be formatted as "Title FirstName LastName"
%   (where Title is "Dr." or "Prof." for example),
% - degree programme  should be "BSc", "MEng", "MSci", "MSc" or "PhD",
% - dissertation title should be correctly capitalised (plus you can have
%   an optional sub-title if appropriate, or leave this field blank),
% - dissertation type should be formatted as one of the following:
%   * for the MEng degree programme either "enterprise" or "research" to
%     reflect the stream,
%   * for the MSc  degree programme "$X/Y/Z$" for a project deemed to be
%     X%, Y% and Z% of type I, II and III.
% - year              should be formatted as a 4-digit year of submission
%   (so 2014 rather than the academic year, say 2013/14 say).

\documentclass[ % the name of the author
                    author={Daniel Page},
                % the name of the supervisor
                supervisor={Dr. Andrew Calway},
                % the degree programme
                    degree={MSc},
                % the dissertation    title (which cannot be blank)
                     title={Some Structural Guidelines for Data Science MSc Theses, Including Those With Long Titles that Run Across Multiple Lines on the Front Page},
                % the dissertation subtitle (which can    be blank)
                  subtitle={And those including an optional subtitle too, for good measure},
                % the dissertation     type
                      type={},
                % the year of submission
                      year={2021}]{dissertation}

\begin{document}

% =============================================================================

% This section simply introduces the structural guidelines.  It can clearly
% be deleted (or commented out) if you use the file as a template for your
% own dissertation: everything following it is in the correct order to use 
% as is.

\section*{Prelude}
\thispagestyle{empty}


A typical MSc thesis will be structured according to a fairly standard sequence of 
sections, and one example is shown here.  However, it is hard and perhaps even 
counter-productive to generalise: the goal of this document is {\em not}\/ to be prescriptive, to show you how you {\em must} structure your thesis; rather it is to simply to act as a guideline, to show you how you {\em could} structure your thesis.  In particular, each chapter's indicated page-count is
important but {\em not}\/ absolute: the aim is simply to highlight that a 
clear and concise description is better than a rambling alternative that makes
it hard to separate important content and facts from irrelevant trivia.

You can use this document as a \LaTeX-based 
template for your own thesis by simply deleting extraneous sections
and content, and adding your own text. If you do that, we recommend using BibTeX
to deal with the associated bibliography. 
You may want to use Overleaf (an online cloud-based collaborative \LaTeX platform).


The words ``thesis'' and ``dissertation'' are used interchangeably here. All active hyperlinks in the PDF file are rendered in the same shade of dark blue. 



% =============================================================================

% This macro creates the standard UoB title page by using information drawn
% from the document class (meaning it is vital you select the correct degree 
% title and so on).


\maketitle

% After the title page (which is a special case in that it is not numbered)
% comes the front matter or preliminaries; this macro signals the start of
% such content, meaning the pages are numbered with Roman numerals.

\frontmatter

% This macro creates the standard UoB declaration; on the printed hard-copy,
% this must be physically signed by the author in the space indicated.

\makedecl

% LaTeX automatically generates a table of contents, plus associated lists 
% of figures, tables and algorithms.  The former is a compulsory part of the
% dissertation, but if you do not require the latter they can be suppressed
% by simply commenting out the associated macro.

\tableofcontents
\listoffigures
\listoftables
\listofalgorithms
\lstlistoflistings

% The following sections are part of the front matter, but are not generated
% automatically by LaTeX; the use of \chapter* means they are not numbered.

% -----------------------------------------------------------------------------

\chapter*{Abstract}

{\bf A compulsory section, of at most $1$ page} 
\vspace{1cm} 

\noindent
This section should summarise the project context, aims and objectives,
and main contributions (e.g., deliverables) and achievements.  The goal is to ensure that the 
reader is clear about what the topic is, what you have done within this 
topic, {\em and}\/ {\bf what your view of the outcome is.}

Essentially 
this section is a (very) short version of what is typically covered in more depth in the first 
chapter.  If appropriate, you should include here  
a clear statement of your research hypothesis.  This will obviously differ significantly
for each project, but an example might be as follows:

\begin{quote}
My research hypothesis is that a suitable genetic algorithm will yield
more accurate results (when applied to the standard ACME data set) than 
the algorithm proposed by Jones and Smith, while also executing in less
time.
\end{quote}

\noindent
The latter aspects should (ideally) be presented as a concise, factual 
list of the main points of achievement.  Again the points will differ for each project, but 
an might be as follows:

\begin{quote}
\noindent
\begin{itemize}
\item I spent $120$ hours collecting material on and learning about the 
      Java garbage-collection sub-system. 
\item I wrote a total of $5000$ lines of {\em Python} source code, and associated orchestration scripts. 
\item I designed a new algorithm for computing the non-linear mapping 
      from A-space to B-space using a genetic algorithm.
\item I implemented a version of the algorithm proposed by Jones and 
      Smith (2010), corrected a mistake in it, and 
      compared the results with several alternatives.
\end{itemize}
\end{quote}

% -----------------------------------------------------------------------------

\chapter*{Summary of Changes}

{\bf A conditional section, of at most $1$ page} 
\vspace{1cm} 

If (and only if) the dissertation represents a resubmission (e.g., as the result of
a resit), then this section is compulsory: the content should summarise all
non-trivial changes made to the initial submission.  Otherwise you can
omit it, since {\bf a summary of this type is only needed for resubmissions}.

When included, the section will ideally be used to highlight additional
work completed, and address criticism raised in any associated feedback.
Clearly it is difficult to give generic advice about how to do so, but
an example might be as follows:

\begin{quote}
\noindent
\begin{itemize}
\item Feedback from the initial submission criticised the design and 
      implementation of my genetic algorithm, stating ``there seems 
      to have been no attention to computational complexity during the
      design, and obvious methods of optimisation are missing within
      the resulting implementation''.  Chapter 3 now includes a
      comprehensive analysis of the algorithm, in terms of both time
      and space.  While I have not altered the algorithm itself, I
      have included a cache mechanism (also detailed in Chapter 3)
      that provides a significant improvement in average run-time.
\item I added a feature in my implementation to allow automatic rather
      than manual selection of various parameters; the experimental
      results in Chapter 4 have been updated to reflect this.
\item Questions after the presentation highlighted a range of related
      work that I had not considered: I have make a number of updates 
      to Chapter 2, resolving this issue.
\end{itemize}
\end{quote}

% -----------------------------------------------------------------------------

\chapter*{Supporting Technologies}

{\bf A compulsory section, of at most $1$ page}
\vspace{1cm} 

\noindent
This section should present a detailed summary, in bullet point form, 
of any third-party resources (e.g., hardware and software components) 
used during the project.  Use of such resources is always perfectly 
acceptable: the goal of this section is simply to be clear about how
and where they are used, so that a clear assessment of your work can
result.  The content can focus on the project topic itself (rather,
for example, than including ``I used \mbox{\LaTeX} to prepare my 
dissertation''); an example is as follows:

\begin{quote}
\noindent
\begin{itemize}

\item I used the {\em Pandas} and {\em Seaborn} public-domian Python Libraries. 

\item I used a parts of the OpenCV computer vision library to capture 
      images from a camera, and for various standard operations (e.g., 
      threshold, edge detection).

\item I used Amazon Web Services for remote storage and processing of data. Specifically, I used:
 \begin{itemize}
 \item Simple Storage Service (S3) for data storage
 \item Elastic Compute Cloud (EC2) for provision of virtual machines
 \item Elastic Beanstalk for scaling and load management
 \item Sagemaker for all the machine learning components of my project. 
 \end{itemize}
 
\item I used \LaTeX\ to format my thesis, via the online service {\em Overleaf}. 
\end{itemize}
\end{quote}

% -----------------------------------------------------------------------------

\chapter*{Notation and Acronyms}

{\bf An optional section, of roughly $1$ or $2$ pages}
\vspace{1cm} 

\noindent
Any well written document will introduce notation and acronyms before
their use, {\em even if} they are standard in some way: this ensures 
any reader can understand the resulting self-contained content.  

Said introduction can exist within the dissertation itself, wherever 
that is appropriate.  For an acronym, this is typically achieved at 
the first point of use via ``Advanced Encryption Standard (AES)'' or 
similar, noting the capitalisation of relevant letters.  However, it 
can be useful to include an additional, dedicated list at the start 
of the dissertation; the advantage of doing so is that you cannot 
mistakenly use an acronym before defining it.  A limited example is 
as follows:

\begin{quote}
\noindent
\begin{tabular}{lcl}
AES                 &:     & Advanced Encryption Standard                                         \\
DES                 &:     & Data Encryption Standard                                             \\
                    &\vdots&                                                                      \\
${\mathcal H}( x )$ &:     & the Hamming weight of $x$                                            \\
${\mathbb  F}_q$    &:     & a finite field with $q$ elements                                     \\
$x_i$               &:     & the $i$-th bit of some binary sequence $x$, st. $x_i \in \{ 0, 1 \}$ \\
\end{tabular}
\end{quote}

% -----------------------------------------------------------------------------

\chapter*{Acknowledgements}

{\bf An optional section, of at most $1$ page}
\vspace{1cm} 

\noindent
It is common practice (although totally optional) to acknowledge any
third-party advice, contribution or influence you have found useful
during your work.  Examples include support from friends or family, 
the input of your Supervisor and/or Advisor, external organisations 
or persons who  have supplied resources of some kind (e.g., funding, 
advice or time), and so on.

\vspace{1cm}
Dave Cliff writes here to say huge thanks to his colleague Dr Dan Page for sharing this \LaTeX\ thesis template, which was originally written by Dan, for Computer Science dissertations. Dave edited Dan's original to better suit the needs of the Data Science MSc: please don't hassle Dan about any of this, but do feel free to contact Dave if you have any questions or comments on it.  

% =============================================================================

% After the front matter comes a number of chapters; under each chapter,
% sections, subsections and even subsubsections are permissible.  The
% pages in this part are numbered with Arabic numerals.  Note that:
%
% - A reference point can be marked using \label{XXX}, and then later
%   referred to via \ref{XXX}; for example Chapter\ref{chap:context}.
% - The chapters are presented here in one file; this can become hard
%   to manage.  An alternative is to save the content in seprate files
%   the use \input{XXX} to import it, which acts like the #include
%   directive in C.

\mainmatter

% -----------------------------------------------------------------------------

\chapter{Introduction: Contextual Background}
\label{chap:introduction}
Chronic pain is a common and complex condition that persists for longer than three months, affecting 30-33\% of the population worldwide and 1 in 3 adults in the UK \cite{shetty_chronic_2025}. A meta-analysis of 19 populations by \cite{fayaz_prevalence_2016} showed that 43.5\% of the UK population (28 million adults) experience chronic pain and 14.3\% (8 million adults) live with moderate or severe disabling chronic pain. Due to the ageing population in England, chronic pain is estimated to grow by 1.7 million by 2024 \cite{toby_watt_health_2023}. Patients with chronic pain can also experience impacts on their quality of life, daily living activities, and psychological well-being \cite{andersen_psychosocial_2014}. The treatment for chronic pain is limited because the underlying cause of this condition is unknown. The standard treatments for acute pain are opioids, but they are not considered safe for long term use for chronic pain due to the shifts in the physiology of the nociceptors. 
Chronic pain is caused by injury or disease in the nervous system and is associated with persistent hyperactivity of nociceptors, which are specialised to detect damaging stimuli and/or inflammation in rodents or humans \cite{walters_persistent_2023}. Therefore, a constant activation of these nociceptors, especially the C-fibre nociceptor, is a driver of prolonged pain and can lead to maladaptive neuroplasticity. Therefore, measuring the abnormal activities of these nociceptors could reveal and allow for precise pain targeting \cite{ossipov_central_2010}. These nociceptor activities can be recorded using microneurography, a technique where an electrode is inserted through the skin into a nerve bundle to record action potentials from a single fibre of interest. However, the spike patterns of these C-fibre nociceptors in the microneurography recordings are challenging because of the small spike amplitude, low signal-to-noise ratio, the conduction velocity and activity dependence, and their waveform shapes are similar between different fibres \cite{troglio_hybrid_2026}. Comprehensive spike detection and sorting analysis are therefore required. Many studies have researched and developed methods for extracellular spike detection and sorting for multi-electrode data, but not many focus on single electrodes where the spatial information is absent, and the discharge rates of unmyelinated peripheral nerve fibres are lower. 

\section{Structure}



% -----------------------------------------------------------------------------

\chapter{Technical Background}
\label{chap:background}
\section{Extracellular Recordings}
\label{sec:extracellular}
\section{Extracellular spike detection}
\label{sec:detection}
\section{Spike sorting and classification}
\label{sec:sorting}
\section{Noise and signal quality of extracellular recordings}
\label{sec:noise}
\section{Signal preprocessing and denoising}
\label{sec:denoising}
\subsection{Wavelet-based denoising}
\subsection{Deep learning-based denoising}
\section{Research gap}
\label{sec:research_gap}

% -----------------------------------------------------------------------------

\chapter{Execution}
\label{chap:execution}
This chapter describes the data and the methods used to detect, classify and denoise the target fibre's spikes in a single channel. It first introduces the recording and how the target unit was identified, then estimates the noise floor. It then describes the detection pipeline in three stages: building and matching a spike template, and constraining detections by their timing after the stimulus. It then describes the two classifiers used on the detection survivors: a one-class and two-class support vector machine. Finally describes three denoising methods: Wavelet, Dense Autoencoder and Convolutional Autoencoder to evaluate each method and to asses whether denoising improves detection or classification, and the metrics used throughout. 
\section{The data: rat saphenous nerve microneurography}
This project uses an extracellular microneurography recording from the peripheral nerve of a rat, sampled at 30 kHz over approximately 2.3 hours. The nerve carries many afferent fibres, so the electrode records the action potentials of many units at once, superimposed on background noise. Recordings were made using two multi-electrode sensor units placed at different positions along the nerve, each containing several
channels corresponding to slightly different electrode locations. Eight channels carry manually labelled ground truth spikes (CH3, CH8, CH9, CH10, CH14, CH19, CH25, CH30). During the recording, an electrical stimulus (Pain stimulus) was delivered repeatedly, giving 2,310 stimulation pulses (TTL) at a normal rate of around 0.25 Hz. The target fibre responds to each stimulus by firing at an approximately fixed latency. The ground truth spike times were marked manually by a human annotator \ref{fig:raw signal}, who only
labelled recognisable nerve responses; therefore, the annotations are conservative, and not every response is labelled. Each spike is marked at the point where the signal crosses zero on its descending slope. 
\begin{figure}[htbp]
    \centering
    \includegraphics[width=\textwidth]{figures/raw signal with target spikes in blue and stimuli dashed red.png}
    \caption{Segment representing the raw signal of channel 19 with the manually annotated spikes in blue and the stimulus times in red. Each target spike follows a stimulus at a latency of around 100 ms.}
    \label{fig:raw signal}
\end{figure}
\subsection{Target unit and channels}
The lab identified a single target unit as a strong unit with large identifiable spikes on multiple channels. This project analyses five of these channels: 9, 14, 19, 25, and 30, with a strong signal, and each is analysed independently. 
The spike latencies of the target unit were manually annotated on each channel; consequently, they are treated as reference annotations rather than complete ground truth, as they were manually identified and may contain errors. 
Only a 15-minute window was used for analyses, beginning at the 900th stimulus pulse. 
This window length was chosen because the target unit's latency remained stable; a longer window showed activity-dependent slowing that spreads the latency and prevents a fixed latency window from capturing the spikes. The signals were preprocessed by the recording laboratory, so it has been band-pass filtered 300-6000 Hz, common average referenced, and stimulus blanked. 
\section{Noise estimation and preprocessing}
As mentioned, no further filtering was applied in this project, because the recording given was already bandpass-filtered, with stimulation artefact blanked (set to zero). 
The target units are small, only about two to three times the amplitude of the background noise; consequently, they do not stand out clearly and can not be isolated by amplitude threshold alone. Therefore, removing the noise, detecting and classifying the target unit's spikes are the central challenges this chapter tries to address. The noise level is estimated using the median absolute deviation (MAD) \cite{quiroga_unsupervised_2004, park_deep_2019}:
\begin{equation}
    \sigma_{\text{noise}} = \frac{\operatorname{median}(|x|)}{0.6745}
\end{equation}
where $x$ is the signal. The median is used instead of standard deviation because the spikes are infrequent, so the median does not inflate the estimate the way standard deviation would; dividing by 0.6745 rescales the value to be comparable to a standard deviation under an assumption of Gaussian noise \cite{drebitz_optimizing_2019}. MAD is an estimate, so it assumes the signal is mostly noise and approximately Gaussian \cite{coupe_robust_2010}. But because it is applied identically to every signal, it is a valid basis for comparison.
\section{Leakage safe evaluation protocol}
For unbiased performance, a chronological split was used on the recording by dividing the recording into a development set on the first 70\% [3602.7-4232.7 s] and a held-out evaluation on the last 30\% of the recording [4232.7-4502.7 s]. All hyperparameter tuning was conducted on the development set, and the evaluation set was scored once to prevent data leakage and to avoid overfitting \cite{rashid_spike_2025}. All component that could carry information across the split was fitted on the development set only; the spike template was built from target spikes in the development annotations split only; the autoencoders were trained on synthetic data derived from the development set. Detection parameters like shape similarity and latency window were selected by maximising F1 on the development set. The classifiers and the two-class classifier were fitted on the development survivors only with fixed hyperparameters and scored once on the evaluation survivors. The confidence intervals were computed by stratified bootstrap resampling of the evaluation candidates.  

\section{Detection}
\subsection{Template matching and spike detection}
The spike template was built by averaging the annotated unit's spikes on the development set, to prevent leakage in later stages. For each marked spike, a 90-sample window (45 samples on either side of the spike trough) is then extracted from the signal and then averaged to give the unit's mean waveform. 
\section{Denoising methods}
\section{Classification}
\section{Evaluation metrics}
\section{Multi-channel consistency}

{\bf A topic-specific chapter, roughly 30\% of the total page-count}
\vspace{1cm} 

\noindent
This chapter is intended to describe what you did: the goal is to explain
the main activity or activities, of any type, which constituted your work 
during the project.  The content is highly topic-specific. For some 
projects it will make sense to split the content into two main sections, or maybe even into two separate chapters: one 
will discuss the design of something, including any rationale or decisions made, 
and the other will discuss how this design was realised via some form of 
implementation.  You could instead give this chapter the title ``Design and Implementation"; or you might split this content into two chapters, one titled ``Design" and the other ``Implementation" .

Note that it is common to include evidence of ``best practice'' project 
management (e.g., use of version control, choice of programming language 
and so on).  Rather than simply a rote list, make sure any such content 
is useful and/or informative in some way: for example, if there was a 
decision to be made then explain the trade-offs and implications 
involved.


% -----------------------------------------------------------------------------

\chapter{Critical Evaluation}
\label{chap:evaluation}
\section{Data and channels}
The evaluation was performed on the five channels where the target unit has a strong signal (channels 9, 14, 19, 25, and 30). Each channel is analysed independently as a single-channel recording over 15 minutes window for consistency analysis. This table \ref{tab:channel data overview} shows how many target annotations were in each channel and how many were in the 15-minute window analysed; each channel was split chronologically into a development set used for all tuning and a held-out set that was scored once. The noise floor MAD is consistent between channels. 
\begin{table}[htbp]
    \centering
    \begin{tabular}{lccccc}
    Channel & Total annotations & Window annotations & Developemnt & Evaluation & MAD \\
    \toprule
    9 & 1628 & 170 & 108 & 62 & 4.80 \\
    14 & 1639 & 173 & 113 & 60 & 4.88 \\
    19 & 1623 & 177 & 116 & 61 & 4.96 \\
    25 & 1777 & 186 & 120 & 66 & 4.63 \\
    30 & 1414 & 174 & 133 & 61 & 4.82 \\
    \bottomrule
    \end{tabular}
    \caption{Overview of the five channels analysed for this project. The target unit appears with a strong signal on all channels. Each channel is analysed independently over a 15-minute window with a chronological split of 70\% development and 30\% evaluation. MAD is the raw noise level estimate for each channel.}
    \label{tab:channel data overview}
\end{table}
\section{Denoising quality}
The three denoisers were assessed across the five channels using two measures: MAD to assess each denoiser's noise reduction relative to the raw,  and the preservation of spike amplitude by using peak-to-peak amplitude relative to the raw signal. These measures were used instead of reconstructive error because no clean signal exists for these real recordings to compare against \cite{baldazzi_systematic_2020, adam_comparative_2025}
All denoisers reduced the noise level, as seen in Table \ref{tab:denoising quality}. Dense reduces the noise the most to MDA ~0.50, and the wavelet denoiser the least to ~2.30 compared to raw ~ 4.80. Despite the Dense autoencoder achieving the lowest MAD, it only retained 17-20\% of the raw spike amplitude across channels, making it the most destructive. The convolutional autoencoder reduced the noise and preserved the spike amplitude, achieving 55-57\% \ref{tab:denoising quality} \ref{fig:denosiing quality ch19}. This pattern was consistent across the five channels. 

\begin{table}[htbp]
    \centering
    \begin{tabular}{llccc}
    \toprule
    Channel & Signal & MAD & Peak-to-peak & Retained \\
    \midrule
    \multirow{4}{*}{9} & Raw & 4.80 & 22.99 & 100\% \\
    & Wavelet & 2.23 & 7.36 & 31\% \\
    & Dense & 0.48 & 4.47 & 19\% \\
    & Conv & 1.08 & 13.13 & 55\% \\
    \midrule
\multirow{4}{*}{14} & Raw & 4.89 & 23.82 & 100\% \\
    & Wavelet & 2.23 & 7.36 & 31\% \\
    & Dense & 0.48 & 4.47 & 19\% \\
    & Conv & 1.08 & 13.13 & 55\% \\
\midrule
\multirow{4}{*}{19} & Raw & 4.96 & 24.08 & 100\% \\
    & Wavelet & 2.25 & 7.62 & 32\% \\
    & Dense & 0.58 & 4.81 & 20\% \\
    & Conv & 1.25 & 13.67 & 57\% \\
\midrule
\multirow{4}{*}{25} & Raw & 4.63 & 21.75 & 100\% \\
    & Wavelet & 2.31 & 7.77 & 36\% \\
    & Dense & 0.50 & 3.91 & 18\% \\
    & Conv & 1.05 & 11.56 & 53\% \\
\midrule
\multirow{4}{*}{30} & Raw & 4.82 & 22.75 & 100\% \\
    & Wavelet & 2.25 & 7.22 & 32\% \\
    & Dense & 0.48 & 3.92 & 17\% \\
    & Conv & 1.23 & 12.68 & 56\% \\
    \bottomrule
    \end{tabular}
    \caption{Denoising quality across the five channels. Each denoiser reduces the noise floor (MAD) and preserves the spike amplitude peak-to-peak relative to raw. The dense autoencoder shows the lowest MAD and retains the least amplitude across all channels, while the convolutional autoencoder reduces the noise floor but preserves the most spike amplitude across all channels.}
    \label{tab:denoising quality}
\end{table}
\begin{figure}[htbp]
    \centering
    \includegraphics[width=\textwidth]{figures/denoising_quality_leakage_safe.png}
    \caption{Effect of each denoiser on the signal on channel 19. Left is an overview of a segment around a target spike. Right is a detail view of the spike.}
    \label{fig:denosiing quality ch19}
\end{figure}

This suggests that MAD alone is not enough to determine denoising quality, because a distorted spike carries less information for detection and classification. So, only focusing on noise removal ignores whether the spike morphology is preserved \cite{baldazzi_systematic_2020}. However, \cite{baldazzi_systematic_2020} used RMSE to quantify morphology preservation against a clean reference. Here, the recording is real and noisy with annotated ground truth, so peak-to-peak was measured following \cite{adam_comparative_2025}. The results for denoising quality are consistent across the five channels; the raw and convolutional spikes retain their depth, while the wavelet and dense did not. 

\section{Template matching detection}
This section compares detection before and after denoising, comparing the three denoisers across the five channels. 
\subsection{Baseline detection}
 Detection was applied to the raw signal on each channel and scored against the reference annotations in the held-out evaluation period \ref{tab:Baseline detection}, \ref{fig:detection baseline}. The false positive to true positives ratio (FP/TP across all channels exceeded 1, ranging from 2.75 to 4.73, and recall was low, around 0.18 to 0.27. Detection on the raw signal of every channel showed more false positives than true positives. This indicates that these false positives resemble the target spikes in both waveforms and latency window, as they passed both the shape similarity threshold and the latency window. These results are consistent with the literature. According to \cite{troglio_hybrid_2026, hill_quality_2011} in single-electrode microneurography analysis, waveform shapes are highly similar across fibres, and it is difficult to separate overlapping units on a single electrode. The TP/FP ratio here for the raw signal is around 2.75 to 4.73, which is lower than those reported by \cite{troglio_hybrid_2026} for multi-fibre recordings; the FP/TP is around 2 to 20. However, these studies did not denoise the signal; the next subsection evaluates whether denoising the signal would make spike detection by shape and latency any better. 

 \begin{table}[htbp]
     \centering
     \begin{tabular}{lcccccccc}
     \toprule
    Channel & Survivors & TP & FP  & FN & Precision & Recall & F1 & FP/TP \\
    \midrule
    9  & 61 & 12 & 49 & 50 & 0.197 & 0.194 & 0.195 & 4.08 \\
    14 & 60 & 16 & 44 & 44 & 0.267 & 0.267 & 0.267 & 2.75 \\
    19 & 63 & 11 & 52 & 50 & 0.175 & 0.180 & 0.177 & 4.73 \\
    25 & 58 & 14 & 44 & 52 & 0.241 & 0.212 & 0.226 & 3.14 \\
    30 & 62 & 13 & 49 & 48 & 0.210 & 0.213 & 0.211 & 3.77 \\
    \bottomrule 
     \end{tabular}
     \caption{Per-channel baseline detection on raw signal on held-out evaluation. The false positives to true positive ratio (FP/TP exceeds 1 for every channel, with channel 19 the highest and channel 14 the lowest.}
     \label{tab:Baseline detection}
 \end{table}
\begin{figure}[H]
     \centering
     \includegraphics[width=\textwidth]{figures/Detection performamce befoer denoiisng (1).png}
     \caption{Baseline detection performance per channel on the raw signal. Spike detection performance using similarity thresholding and a latency window: TP, FP, and FN are visualised on a log scale for the raw signal across the five channels on the held-out evaluation period. All channels have higher false positives compared to the true positives, with Ch19 exhibiting the most false positives and the least true positives.}
     \label{fig:detection baseline}
 \end{figure}
\subsection{Detection after denoising}
Detection was then applied to each denoised signal from wavelet, dense and convolutional and compared to the raw baseline \ref{tab:detection denoised}, \ref{fig:detection after denoising}. 
The effect of denoising on the FP/TP is inconsistent across channels, and no denoiser improved the FP/TP; the ratio remains above 1 for all denoisers and all channels. 
Wavelet improved the ratio slightly on three channels (ch9, ch14, and ch19), and worsened the ratio on two channels (ch25 and ch30); dense also improved three (ch9, ch25, and ch30) and worsed the ratio on two channels (ch14, 19); convlutional improved 2 (ch9, ch19) and showed worse ratio on (ch14, ch25, ch30). Only channel 9: the three denoisers reduced FP/TP, although even there the ratio remained above 2.7.

\begin{table}[htbp]
    \centering
    \begin{tabular}{llccccccc}
    \toprule
    Channel & Signal & Survivors & TP & FP & FN & Recall & F1 & FP/TP \\
    \midrule
    \multirow{4}{*}{9}  
    & Raw & 61 & 12 & 49 & 50 & 0.194 & 0.195 & 4.08 \\
    & Wavelet & 59 & 16 & 43 & 46 & 0.258 & 0.264 & 2.69 \\
    & Dense & 60 & 15 & 45 & 47 & 0.242 & 0.246 & 3.00 \\
    & Conv & 65 & 16 & 49 & 46 & 0.258 & 0.252 & 3.06 \\
    \midrule
    \multirow{4}{*}{14} 
    & Raw & 60 & 16 & 44 & 44 & 0.267 & 0.267 & 2.75 \\
    & Wavelet & 59 & 16 & 43 & 44 & 0.267 & 0.269 & 2.69 \\
    & Dense & 55 & 10 & 45 & 50 & 0.167 & 0.174 & 4.50 \\
    & Conv & 65 & 15 & 50 & 45 & 0.250 & 0.240 & 3.33 \\
    \midrule
    \multirow{4}{*}{19} 
    & Raw & 63 & 11 & 52 & 50 & 0.180 & 0.177 & 4.73 \\
    & Wavelet & 58 & 14 & 44 & 47 & 0.230 & 0.235 & 3.14 \\
    & Dense & 64 & 8 & 56 & 53 & 0.131 & 0.128 & 7.00 \\
    & Conv & 63 & 16 & 47 & 45 & 0.262 & 0.258 & 2.94 \\
    \midrule
    \multirow{4}{*}{25} 
    & Raw & 58 & 14 & 44 & 52 & 0.212 & 0.226 & 3.14 \\
    & Wavelet & 52 & 12 & 40 & 54 & 0.182 & 0.203 & 3.33 \\
    & Dense & 61 & 21 & 40 & 45 & 0.318 & 0.331 & 1.90 \\
    & Conv & 61 & 12 & 49 & 54 & 0.182 & 0.189 & 4.08 \\
    \midrule
    \multirow{4}{*}{30} 
    & Raw & 62 & 13 & 49 & 48 & 0.213 & 0.211 & 3.77 \\
    & Wavelet & 51 & 7 & 44 & 54 & 0.115 & 0.125 & 6.29 \\
    & Dense & 65 & 14 & 51 & 47 & 0.230 & 0.222 & 3.64 \\
    & Conv & 62 & 11 & 51 & 50 & 0.180 & 0.179 & 4.64 \\
    \bottomrule
    \end{tabular}
    \caption{Per channel and signal detection after denoising on the held-out evaluation. Denoising did not reduce the FP/TP ratio; all are exceeding 1 on every channel, and it shows an inconsistent effect across channels.}
    \label{tab:detection denoised}
\end{table}
\begin{figure}[htbp]
     \centering
     \includegraphics[width=\textwidth]{figures/Detection performamce after denoisng.png}
     \caption{Detection performance per channel before and after denoising the signal. The TP, FP, and FN are visualised on a log scale for the denoised signal across the five channels on the held-out evaluation period. Denoising does not consistently reduce the false positive count.}
     \label{fig:detection after denoising}
 \end{figure}

When comparing spike template matching detection before denoising to after denoising, none of the three denoisers reduced the FP/TP ratio to below 1 on any channel; even the largest improvement, which was dense on CH25, only improved it to 1.90 compared to the raw ratio of 3.14. The scattered pattern of improvement and worsening is quite balanced across the fifteen denoiser channel combinations. 

Channel-to-channel variation [TO DO]



Several reasons and mechanisms reported in the literature for denoising increasing the number of false positives with some denoisers and some channels, and these fall into two categories. These false positives could be genuine biological spikes that were buried by the noise and were detected by the template matching after denoising, or artefacts introduced by the denoisers themselves.
According to \cite{bengio_generalized_2013}, denoising autoencoders can learn to pull corrupted points back onto the manifold via a learned vector field, so when the model receives unstructured noise, it can reshape that noise into the nearest pattern, which, in this case, is spikes. This could be the reason for the higher number of false positives compared to raw. Furthermore, denoising can diminish the spikes' fine features that help distinguish them from neighbouring units \cite{baldazzi_systematic_2020}, which in turn causes non-target spikes to resample the target template. If the additional detections (false positives) by the denoisers are actual spikes, they should look like the spike's waveforms in the raw signal; and if they are artefacts, they should be absent or near flat in the raw signal since they were created. This check was only done on ch19 for the dense autoencoder, where the FP/TP ratio increased from 4.73 to 7.00 with 11 additional detections. These additional detections have a mean trough of -11.77 in the raw signal, similar to the target spikes, indicating that these extra false positives might be spike-like events rather than artefacts produced by the denoiser. However, this can not rule out any other possibilities, as only the mean trough spikes were checked and only on one channel with one denoiser. Nonetheless, in a single channel, spikes do overlap, have similar waveform shapes and latency, and without the spatial information that of multi-channel analysis, it is difficult to distinguish target spikes from non-target \cite{troglio_hybrid_2026}.  
\section{Classification}

\subsection{Baseline classification}
The one-class and two-class SVMs were applied to the survivors from detection on the raw signal, trained on the development set and scored on the held-out evaluation set once.  Both the table and the figure \ref{tab:classification baseline} \ref{fig:classification_baseline} show the ROC-AUC with 95\% confidence interval included 0.5. For one class svm the ROC-AUC for all channels ranged from 0.455 to 0.602, while the two-class SVM ranged from 0.399 to 0.606, and all confidence intervals included 0.5. The confusion matrix counts \ref{tab:classification baseline confusion metrics} show that for every channel, the false positives are more than true positives and the precision is between 0.15 and 0.30. 
\begin{table}[H]
    \centering
    \begin{tabular}{lcc}
    \toprule
    Channel &  One-class ROC-AUC [95\% CI] & Two-class ROC-AUC [95\% CI] \\
    \midrule
    9 & 0.602 [0.446, 0.755] & 0.481 [0.299, 0.658] \\
    14 & 0.477 [0.303, 0.651] & 0.506 [0.347, 0.666] \\
    19 & 0.455 [0.264, 0.652] & 0.502 [0.283, 0.719] \\
    25 & 0.539 [0.370, 0.711] & 0.399 [0.237, 0.571] \\
    30 & 0.523 [0.356, 0.684] & 0.606 [0.411, 0.780] \\
    \bottomrule
    \end{tabular}
    \caption{Per-channel baseline classification performance on the raw signal. All ROC-AUC confidence intervals are wide, but all include 0.5, indicating that the model's performance on the raw signal is indistinguishable from chance.}
    \label{tab:classification baseline}
\end{table}
\begin{table}[H]
    \centering
    \begin{tabular}{llccccccc}
    Classifier & Channel & TP & FP & TN & Precision & Recall & F1 \\
    \midrule
    \multirow{5}{*}{One-class}
    & 9  & 6  & 19 & 6 & 30 & 0.240 & 0.500 & 0.324 \\
    & 14 & 10 & 28 & 6 & 16 & 0.263 & 0.625 & 0.370 \\
    & 19 & 6  & 30 & 5 & 22 & 0.167 & 0.545 & 0.255 \\
    & 25 & 8  & 19 & 6 & 25 & 0.296 & 0.571 & 0.390 \\
     & 30 & 8  & 28 & 5 & 21 & 0.222 & 0.615 & 0.327 \\
    \midrule
    \multirow{5}{*}{Two-class}
     & 9  & 7  & 28 & 5 & 21 & 0.200 & 0.583 & 0.298 \\
     & 14 & 12 & 33 & 4 & 11 & 0.267 & 0.750 & 0.393 \\
     & 19 & 6  & 33 & 5 & 19 & 0.154 & 0.545 & 0.240 \\
     & 25 & 7  & 26 & 7 & 18 & 0.212 & 0.500 & 0.298 \\
     & 30 & 8  & 22 & 5 & 27 & 0.267 & 0.615 & 0.372 \\
    \bottomrule
    \end{tabular}
    \caption{Per-channel confusion matrix counts of baseline classification. The true positives and precision scores exceed true positives.}
    \label{tab:classification baseline confusion metrics}
\end{table}
\begin{figure}[H]
    \centering
    \includegraphics[width=\textwidth]{figures/classification_baseline.png}
    \caption{Dot plot showing per-channel ROC-AUC scores for baseline classification on raw signal. The ROC-AUC is shown with 95\% bootstrap confidence intervals for one class on the left side and two classes on the right side. Every interval includes the chance level of 0.5 where the red dashed line.}
    \label{fig:classification_baseline}
\end{figure}
\subsection{Classification after denoising}
On denoised signals, no consistent improvement in classification performance was observed across channels \ref{tab:classificatio-roc denoised}, \ref{fig:classification roc after denoising}. The biggest improvement observed with one-class SVM after denoising the signal was for ch30, where ROC-AUC increased from 0.523 (95\% CI of 0.356-0.684) to 0.675 (95\% CI of 0.409-0.896) following wavelet denoising and 0.649 (95\% CI of 0.483-0.791) following conv denoising. An improvement was also observed on ch14 with dense denoising, where ROC-AUC increased from 0.477 (95\% CI of 0.303-0.651) to 0.589(95\% CI of 0.411-0.764). In contrast, the ch19 denoising signal showed decreased performance, with ROC-AUC decreasing from 0.455 (95\% CI of 0.264-0.652) to 0.370(95\% CI of 0.215-0.522) with the conv denoiser. on the other hand, wavelet reduced two-class performance for ch30 from ROC-AUC 0.606(95\% CI of 0.411-0.780) to 0.344(95\% CI of 0.133-0.578). The full confusion matrix after denoising is reposted in the Appendix \ref{}. 
\begin{table}[H]
    \centering
    \begin{tabular}{llcc}
    \toprule
    Channel & Denoiser & One-class ROC-AUC [95\% CI] & Two-class ROC-AUC [95\% CI] \\
    \midrule
    \multirow{3}{*}{9}  
    & Wavelet & 0.541 [0.366, 0.711] & 0.509 [0.331, 0.676] \\
    & Dense & 0.536 [0.370, 0.702] & 0.527 [0.351, 0.704] \\
    & Conv & 0.537 [0.363, 0.703] & 0.509 [0.356, 0.672] \\
    \midrule
    \multirow{3}{*}{14} 
    & Wavelet & 0.461 [0.305, 0.625] & 0.491 [0.308, 0.673] \\
    & Dense & 0.589 [0.411, 0.764] & 0.613 [0.387, 0.822] \\
    & Conv & 0.517 [0.367, 0.680] & 0.585 [0.419, 0.740] \\
    \midrule
    \multirow{3}{*}{19} 
    & Wavelet & 0.435 [0.258, 0.612] & 0.484 [0.292, 0.680] \\
    & Dense & 0.438 [0.239, 0.623] & 0.531 [0.310, 0.757] \\
    & Conv & 0.370 [0.215, 0.522] & 0.387 [0.243, 0.542] \\
    \midrule
    \multirow{3}{*}{25} 
    & Wavelet & 0.552 [0.371, 0.727] & 0.348 [0.192, 0.521] \\
    & Dense & 0.576 [0.421, 0.720] & 0.380 [0.240, 0.523] \\
    & Conv & 0.580 [0.427, 0.735] & 0.587 [0.422, 0.747] \\
    \midrule
    \multirow{3}{*}{30} 
    & Wavelet & 0.675 [0.409, 0.896] & 0.344 [0.133, 0.578] \\
    & Dense & 0.598 [0.420, 0.758] & 0.410 [0.267, 0.564] \\
    & Conv & 0.649 [0.483, 0.791] & 0.613 [0.415, 0.795] \\
    \bottomrule
    \end{tabular}
    \caption {Per-channel classification performance on denoised signal. The ROC-AUC confidence interval is wide, but all include 0.5; denoised signal does not improve classification to above chance.}
    \label{tab:classificatio-roc denoised}
\end{table}
\begin{figure}[H]
    \centering
    \includegraphics[width=\textwidth]{figures/classification after denoising.png}
    \caption{Per-channel and denoiser (wavelet, dense, and convolutional autoencoders) ROC-AUC score for classification performance on denoised signal.The ROC-AUC is shown with 95\% bootstrap confidence intervals for one class on the left side and two classes on the right side. Denoising did not raise performance above chance; every interval includes the chance level of 0.5, where the red dashed line}
    \label{fig:classification roc after denoising}
\end{figure}
Across all denoisers and channels, the confidence intervals are wide, spanning from 0.2 to 0.8. These intervals were computed using stratified bootstrap resampling, so this width is not due to shifting class balance across resamples. However, two factors account for this wide interval: first, the small evaluation set, which was around 60 survivors per channel, of which only around 14 to 22 are target, limiting the statistical power and widening the intervals. The second factor is that the target and non-target waveforms overlap, making it difficult for the classifier to separate them even on a denoised signal, because non-target events appear to have a spike shape similar to the target \cite{troglio_hybrid_2026, hill_quality_2011}. The confusion matrix \ref{tab:confusion} supports this; the classifier accepts non-target survivors as targets, returning precision [0.13 to 0.40]; false positives exceed true positives, which confirms that classification is not discriminating between target spike and non-target spike and that this is consistent between channels. 
These results suggest that the raw waveform features and the SVM tested here are unable to separate target units from non-target spikes on a single channel. This is a negative finding, suggesting that waveform shape alone can not resolve this contamination. Denoising the signal does not change this conclusion. Although some point estimates differ slightly, none of these differences seems meaningful. For example, the two-class ROC-AUC for the conv denoiser reached 0.61 on ch14 and 30; their confidence intervals still include 0.5 (0.42-0.82) and (0.42-0.80); does not indicate above-chance performance. Similarly, the channel-to-channel variation observed is expected because of the small evaluation set, resulting in ROC-AUC estimates that are quite uncertain, so a difference such as 0.44 vs 0.68 could just be sampling variability rather than evidence that some channels provide better class separation. Furthermore, because multiple channels, denoising and classification conditions are being evaluated, some ROC-AUC may rise due to sampling variability and should therefore be interpreted with caution \cite{sullivan_facts_2021}. 
\section{Results Summary}

% -----------------------------------------------------------------------------

\section{Limitations}
Several limitations must be considered when interpreting the results. 
\chapter{Conclusion}
\section{Future work}
\label{chap:conclusion}

{\bf A compulsory chapter,  roughly 10\% of the total page-count}
\vspace{1cm} 

\noindent
The concluding chapter(s) of a dissertation are often underutilized because they're 
too often left too close to the deadline: it is important to allocate enough time and 
attention to closing off the story, the narrative, of your thesis.

Again, there is no single correct way of closing a thesis. 

One good way of doing this is to have a single chapter consisting of three parts:

\begin{enumerate}
\item (Re)summarise the main contributions and achievements, in essence
      summing up the content.
\item Clearly state the current project status (e.g., ``X is working, Y 
      is not'') and evaluate what has been achieved with respect to the 
      initial aims and objectives (e.g., ``I completed aim X outlined 
      previously, the evidence for this is within Chapter Y'').  There 
      is no problem including aims which were not completed, but it is 
      important to evaluate and/or justify why this is the case.
\item Outline any open problems or future plans.  Rather than treat this
      only as an exercise in what you {\em could} have done given more 
      time, try to focus on any unexplored options or interesting outcomes
      (e.g., ``my experiment for X gave counter-intuitive results, this 
      could be because Y and would form an interesting area for further 
      study'' or ``users found feature Z of my software difficult to use,
      which is obvious in hindsight but not during at design stage; to 
      resolve this, I could clearly apply the technique of Bloggs {\em et al.}.
\end{enumerate}

Alternatively, you might want to divide this content into two chapters: a penultimate chapter with a title such as ``Further Work" and then a final chapter ``Conclusions". Again, there is no hard and fast rule, we trust you to make the right decision. 

And this, the final paragraph of this thesis template, is just a bunch of citations, added to show how to generate a BibTeX bibliography. Sources that have been randomly chosen to be cited here include:
\cite{miller_etal_2018_clojure,webber_marwan_2015,touretzky_2013_lisp,eckmann_etal_1987,marwan_2011,vach_2015,shiller_2017,vytelingum_2006,tesfatsion_2002,rust_etal_1992}.




% =============================================================================

% Finally, after the main matter, the back matter is specified.  This is
% typically populated with just the bibliography.  LaTeX deals with these
% in one of two ways, namely
%
% - inline, which roughly means the author specifies entries using the 
%   \bibitem macro and typesets them manually, or
% - using BiBTeX, which means entries are contained in a separate file
%   (which is essentially a database) then imported; this is the 
%   approach used below, with the databased being dissertation.bib.
%
% Either way, the each entry has a key (or identifier) which can be used
% in the main matter to cite it, e.g., \cite{X}, \cite[Chapter 2}{Y}.

\backmatter

\bibliography{references.bib}

% -----------------------------------------------------------------------------

% The dissertation concludes with a set of (optional) appendices; these are 
% the same as chapters in a sense, but once signalled as being appendices via
% the associated macro, LaTeX manages them appropriately.

\appendix

\chapter{An Example Appendix}
\label{appx:example}

Content which is not central to, but may enhance the dissertation can be 
included in one or more appendices; examples include, but are not limited
to

\begin{itemize}
\item lengthy mathematical proofs, numerical or graphical results which 
      are summarised in the main body,
\item sample or example calculations, 
      and
\item results of user studies or questionnaires.
\end{itemize}

\noindent
Note that in line with most research conferences, the examiners are not
obliged to read such appendices.

% =============================================================================

\end{document}
